\documentclass[11pt]{article}
\usepackage[a4paper,margin=2.45cm]{geometry}
\usepackage[T1]{fontenc}
\usepackage[utf8]{inputenc}
\usepackage{lmodern}
\usepackage{amsmath,amssymb,amsthm,mathtools,bm}
\usepackage{booktabs}
\usepackage{float}
\usepackage{microtype}
\usepackage{graphicx}
\usepackage{xcolor}
\usepackage[hidelinks]{hyperref}
\usepackage[nameinlink,noabbrev]{cleveref}

\definecolor{linkblue}{RGB}{25,75,140}
\hypersetup{colorlinks=true,linkcolor=linkblue,citecolor=linkblue,urlcolor=linkblue}

\newtheorem{theorem}{Theorem}
\newtheorem{proposition}[theorem]{Proposition}
\newtheorem{corollary}[theorem]{Corollary}
\newtheorem{remark}[theorem]{Remark}
\newcommand{\Smat}{\mathbb{S}}
\newcommand{\PSD}{\Smat_+}
\newcommand{\tr}{\operatorname{tr}}
\newcommand{\ip}[2]{\left\langle #1,#2\right\rangle}

\title{Finite-Sample Spectral-Edge Exclusion\\
from Gaussian Correlator Sketches}
\author{Lluís Eriksson\\
\small Independent researcher\\
\small \url{https://github.com/lluiseriksson/finite-sample-spectral-certificates}}
\date{August 2026}

\begin{document}
\maketitle

\begin{abstract}
Spectral gaps are often inferred from noisy Euclidean correlators by fitting a
small number of exponentials or solving a generalized eigenvalue
problem~\cite{Blossier2009}.  Such
point estimates do not by themselves control the error incurred when the
spectral filter is selected from the same data.  We give a finite-sample test
that instead asks whether an entire simultaneous confidence set is compatible
with a proposed upper edge of the visible transfer spectrum.  For Gaussian
sketches of a block correlator, an extreme-singular-value inequality yields a
Loewner confidence band for the unknown joint Gram matrix.  Intersecting this
band with a truncated Hausdorff localizer is one semidefinite feasibility
problem.  We derive its semidefinite alternative: every rejection is accompanied
by a covariance-aware matrix witness, a negative separation margin, and
checkable primal-independent residuals.  The test has type-I error at most
$\alpha$ at fixed truncation order and requires neither known covariance nor a
Gaussian approximation to the quadratic moment features.  A second theorem
gives an explicit witness-dependent sample size for power at least
$1-\alpha$, and one rejection is replayed with exact rational arithmetic.
Simulations in the
interacting axial next-nearest-neighbour Ising chain compare the certificate to
covariance-known ellipsoids, plug-in eigenvectors, fixed Chebyshev filters, and
Ritz estimates.  A 16,800-SDP production grid resolves sample-size and moment-
degree transitions.  We also document two negative results: Hotelling
calibration undercovers the non-Gaussian quadratic features, and the Gaussian
Loewner band undercovers variance-matched Student-$t_5$ primitives.  The method
certifies exclusion of overly large \emph{visible} gaps; it
does not turn non-rejection into a lower bound on a Hamiltonian mass gap.
\end{abstract}

\section{Introduction}

Let $H$ be a Hamiltonian with ground energy $E_0$, choose a Euclidean step
$\tau>0$, and write $T=\exp[-\tau(H-E_0)]$.  Matrix-valued correlator moments
have the form
\begin{equation}
  B_k=\Psi^*T^k\Psi,\qquad k=0,1,\ldots,
  \label{eq:moments}
\end{equation}
where the columns of $\Psi$ are ground-state-subtracted probe states.  The
largest point in the support visible to these probes is
$\rho_{\rm vis}=e^{-\tau\Delta_{\rm vis}}$.  If a probe misses the first
excitation then $\Delta_{\rm vis}$ may exceed the Hamiltonian gap; this
visibility qualification is unavoidable.

Recent work has developed semidefinite bootstraps for exact Euclidean
correlators~\cite{Cho2025}, convergent lower-gap certificates for qubit
Hamiltonians~\cite{Rao2025}, and convergent certified upper bounds on the bulk
gap~\cite{Xu2026}.  Our question is complementary:
\begin{quote}
\emph{Given finitely many correlated random observations of a block
correlator, can one reject a proposed visible gap with simultaneous,
finite-sample error control while choosing the separating filter adaptively?}
\end{quote}
The distinction matters.  A negative eigenvalue of a noisy localizer is a
data-dependent observation, not yet a confidence statement.  Entrywise error
bars ignore correlation; a filter trained on the same sample invalidates a
pointwise interval; and a Gaussian approximation to quadratic observables may
be inaccurate precisely in the moderate-$n$ regime where certification is most
valuable.

Our contribution is a statistical-conic layer that can be placed in front of a
truncated moment localizer.  In the Gaussian-sketch model the primitive sketch
vector is Gaussian, but the reported block moments are quadratic and
Wishart-like.  We construct a confidence set for the primitive covariance
itself, then test its compatibility with the spectral null.  The resulting test
is nonasymptotic and covariance-free.  Its dual supplies an explicit rejection
witness rather than asking the reader to trust an ``infeasible'' solver label.

\paragraph{Scope.}
At finite moment degree, localizer positivity is a necessary condition for a
proposed upper spectral edge.  Its violation rigorously excludes that edge, but
non-violation does not prove a gap.  Consequently our output is a certified
upper bound on a visible gap (or a rejection of an overly large claimed gap),
not a proof of a positive mass gap.  This boundary also distinguishes the
present result from Hamiltonian-algebra hierarchies that can certify lower
bounds~\cite{Rao2025}.

\section{Block moments and the upper-edge localizer}

Assume for clarity that $H$ and the probes admit a real representation; the
complex proper-Gaussian case has different fourth-moment constants and is not
claimed here.  The spectral theorem gives a positive-semidefinite matrix-valued
measure $\mu$ on $[0,1]$ such that
\begin{equation}
  B_k=\int_{[0,1]}x^k\,d\mu(x).
\end{equation}
For a proposed edge $\theta\in(0,1]$ and degree $N$, define
\begin{equation}
  \mathcal L_{N,\theta}(B)
  =\big[\theta B_{i+j}-B_{i+j+1}\big]_{i,j=0}^{N}.
  \label{eq:localizer}
\end{equation}
If $\operatorname{supp}\mu\subseteq[0,\theta]$, then for every vector
polynomial $p(x)=\sum_{i=0}^N x^iv_i$,
\begin{equation}
  \sum_{i,j}v_i^*(\theta B_{i+j}-B_{i+j+1})v_j
  =\int(\theta-x)\,p(x)^*d\mu(x)p(x)\ge0.
\end{equation}
Thus $\mathcal L_{N,\theta}(B)\succeq0$ is a necessary finite-window null.
A violation says that the probes see spectral weight above $\theta$, hence
\begin{equation}
  \Delta_{\rm vis}< -\tau^{-1}\log\theta.
\end{equation}

This elementary localizer is not presented as novel; truncated moment and
localization techniques have a long history~\cite{HeltonNie2011,deKlerk2018}.
The new issue below is how to test it after observing a correlated random Gram
matrix.

\section{Gaussian sketches and a covariance confidence band}

Put $u_{k,a}=T^{k/2}\psi_a$ for time index $k=0,\ldots,2N+1$ and channel
$a=1,\ldots,c$.  Draw independent standard real Gaussian vectors $g_r$ and
observe
\begin{equation}
  z_{r,k,a}=\ip{g_r}{u_{k,a}}.
  \label{eq:sketch}
\end{equation}
After stacking $(k,a)$ into one index, $z_r\sim\mathcal N(0,K)$ with
\begin{equation}
  K_{(k,a),(\ell,b)}
   =\ip{u_{k,a}}{u_{\ell,b}}
   = B_{(k+\ell)/2}[a,b].
  \label{eq:jointgram}
\end{equation}
In particular, the $k$th diagonal time block of $K$ is $B_k$.  Although
$z_{r,k}z_{r,k}^{\mathsf T}$ is an unbiased rank-one estimator of $B_k$, these
quadratic features are not jointly Gaussian.

Let
\begin{equation}
  S=\frac1n\sum_{r=1}^n z_rz_r^{\mathsf T},\qquad
  d=(2N+2)c.
\end{equation}
The following confidence set acts on all times and channels simultaneously.

\begin{theorem}[Wishart--Loewner confidence band]
\label{thm:wishart}
Let $z_1,\ldots,z_n$ be independent $\mathcal N(0,K)$ vectors in
$\mathbb R^d$, where $K\succeq0$ is unknown and the zero mean is known.  For
$0<\alpha<1$, set
\begin{equation}
 \eta=\frac{\sqrt d+\sqrt{2\log(2/\alpha)}}{\sqrt n}<1,
 \quad a=(1+\eta)^{-2},\quad b=(1-\eta)^{-2}.
 \label{eq:bandconstants}
\end{equation}
Then
\begin{equation}
  \mathbb P\{aS\preceq K\preceq bS\}\ge1-\alpha.
  \label{eq:loewnerband}
\end{equation}
The statement remains valid for singular $K$.
\end{theorem}

\begin{proof}
Write the $n\times d$ data matrix as $Z=GK^{1/2}$, with $G$ standard
Gaussian.  The Gaussian extreme-singular-value inequality gives, with failure
probability at most $2e^{-t^2/2}$,
\[
 \sqrt n-\sqrt d-t\le s_{\min}(G)\le s_{\max}(G)
 \le\sqrt n+\sqrt d+t.
\]
With $t=\sqrt{2\log(2/\alpha)}$, division by $\sqrt n$, squaring, and
congruence by $K^{1/2}$ yield
$(1-\eta)^2K\preceq S\preceq(1+\eta)^2K$.  Rearrangement gives
\eqref{eq:loewnerband}.  Congruence does not require invertibility of $K$, so
the same argument covers the singular case.  We use the ambient $d$ rather
than an unknown effective rank.  See~\cite{RudelsonVershynin2010,Vershynin2010}
for the nonasymptotic singular-value bounds.
\end{proof}

\begin{remark}[Known mean is substantive]
If a sample mean is first estimated and subtracted, the degrees of freedom and
constants change.  The present theorem applies to the centered-by-construction
sketches in \eqref{eq:sketch}; it must not be copied unchanged to an unknown-mean
observation model.
\end{remark}

\section{Exact compatibility test and semidefinite alternative}

Let $P(K)=(K_{kk})_{k=0}^{2N+1}$ extract the diagonal time blocks.  Define the
random confidence set
\begin{equation}
  \mathcal C_\alpha(S)=\{Q\in\Smat^d:aS\preceq Q\preceq bS\}.
\end{equation}
We reject the proposed edge $\theta$ precisely when
\begin{equation}
 \text{there is no }Q\in\mathcal C_\alpha(S)
 \text{ with }\mathcal L_{N,\theta}(P(Q))\succeq0.
 \label{eq:primaltest}
\end{equation}
This is one SDP feasibility problem.

\begin{corollary}[Finite-sample validity]
\label{cor:level}
Under the null $\mathcal L_{N,\theta}(P(K))\succeq0$, the probability that
\eqref{eq:primaltest} rejects is at most $\alpha$.
\end{corollary}

\begin{proof}
On the event in \cref{thm:wishart}, the true $K$ is feasible in
\eqref{eq:primaltest}; rejection is therefore impossible.  The rejection event
is contained in the confidence-set failure event.
\end{proof}

No sample splitting is needed: arbitrary adaptive optimization inside the
confidence set preserves this argument.  This is the central statistical
advantage over selecting a polynomial direction from $S$ and attaching an
unadjusted pointwise error bar.

The same event also gives a nonasymptotic power condition.  It makes explicit
how statistical dimension, spectral violation, and sketch count trade off.

\begin{theorem}[Finite-sample detection condition]
\label{thm:power}
Fix an alternative covariance $K\succeq0$ and let
$C_X=(\mathcal L_{N,\theta}\circ P)^*(X)$ for some $X\succeq0$ with
$\tr X=1$.
Suppose
\begin{equation}
 \gamma_X=-\ip{C_X}{K}>0,
 \qquad
 w_X=\tr\left|K^{1/2}C_XK^{1/2}\right|.
 \label{eq:populationmargin}
\end{equation}
If
\begin{equation}
 \delta_+(\eta)w_X<\gamma_X,
 \qquad
 \delta_+(\eta)=\left(\frac{1+\eta}{1-\eta}\right)^2-1,
 \label{eq:powercondition}
\end{equation}
then the test rejects with probability at least $1-\alpha$.  In particular,
writing $r_X=\gamma_X/w_X$ (with the conclusion immediate if $w_X=0$), the
explicit sufficient sample size is
\begin{equation}
 n>
 \left(\sqrt d+\sqrt{2\log(2/\alpha)}\right)^2
 \frac{\big(\sqrt{1+r_X}+1\big)^4}{r_X^2}.
 \label{eq:samplesufficient}
\end{equation}
\end{theorem}

\begin{proof}
On the event used in Theorem~\ref{thm:wishart}, every
$Q\in\mathcal C_\alpha(S)$ obeys
\[
 r_-K\preceq Q\preceq r_+K,
 \qquad
 r_\pm=\left(\frac{1\pm\eta}{1\mp\eta}\right)^2.
\]
Thus, for $E=Q-K$,
$-(1-r_-)K\preceq E\preceq(r_+-1)K$.  Compressing to
$\operatorname{ran}K$ and applying congruence by $K^{\dagger/2}$ shows
\begin{align}
 \ip{C_X}{E}
 &\le (r_+-1)\tr[(K^{1/2}C_XK^{1/2})_+]\\
 &\quad +(1-r_-)\tr[(K^{1/2}C_XK^{1/2})_-]
 \le\delta_+(\eta)w_X.
\end{align}
Condition~\eqref{eq:powercondition} therefore gives
$\ip{X}{\mathcal L_{N,\theta}(P(Q))}=\ip{C_X}{Q}<0$ for every $Q$ in the
confidence set.  No such $Q$ can have a PSD localizer, so rejection occurs on
an event of probability at least $1-\alpha$.  Finally,
$\delta_+(\eta)<r_X$ is equivalent to
$\eta<r_X/(\sqrt{1+r_X}+1)^2$; substituting~\eqref{eq:bandconstants} gives
\eqref{eq:samplesufficient}.
\end{proof}

The bound is deliberately witness-specific and conservative.  It is useful
because it is a genuine power guarantee, not because it predicts the sharp
transition of an optimized SDP.  Optimizing $X$ for empirical power remains
fully protected by Corollary~\ref{cor:level}; using $X$ in
Theorem~\ref{thm:power} as a population design witness is only an analysis
device.

For numerical auditability, let $D$ be any fixed positive diagonal scaling
computed from $S$, put $Q=D\widetilde QD$, and denote the scaled interval by
$A\preceq\widetilde Q\preceq B$.  Let $\mathcal L_D$ be the resulting affine
linear map.

\begin{proposition}[Semidefinite alternative]
\label{prop:dual}
Assume $S\succ0$.  The feasibility problem \eqref{eq:primaltest} is infeasible
if and only if there exist $Y_A,Y_B,Y_L\succeq0$ such that
\begin{align}
 Y_A-Y_B+\mathcal L_D^*(Y_L)&=0,\label{eq:stationarity}\\
 -\ip{Y_A}{A}+\ip{Y_B}{B}&<0.\label{eq:dualmargin}
\end{align}
The multipliers may be normalized by
$\tr Y_A+\tr Y_B+\tr Y_L=1$.
\end{proposition}

\begin{proof}
Pair the three primal inequalities
$\widetilde Q-A\succeq0$, $B-\widetilde Q\succeq0$, and
$\mathcal L_D(\widetilde Q)\succeq0$ with the three PSD multipliers.
Equation~\eqref{eq:stationarity} cancels every term containing
$\widetilde Q$; primal feasibility would then force the constant in
\eqref{eq:dualmargin} to be nonnegative, proving one direction.  Conversely,
the interval is compact and has the strict midpoint $(A+B)/2$ when $S\succ0$.
Strong conic separation therefore supplies the stated multipliers whenever its
intersection with the inverse image of the PSD localizer cone is empty.
Homogeneity permits trace normalization.  If $S$ is singular, the same claim
requires first restricting all matrices to $\operatorname{ran}S$; we do not
silently replace this reduction by a numerical ridge.
\end{proof}

The implementation minimizes the dual margin under trace normalization.  A
reported rejection stores the margin, the Frobenius residual of
\eqref{eq:stationarity}, normalization error, and the smallest eigenvalue of
each numerical multiplier.  The production grid uses a conservative
a-posteriori residual allowance.  In addition, one deliberately wider
confidence interval is rejected by an exact-rational witness whose PSD and
strict-negativity checks use no numerical linear-algebra library; see
Section~\ref{sec:rational}.

\section{Why covariance awareness changes the witness}

The noiseless localizer test uses a minimum eigenvector.  Under simultaneous
uncertainty, that direction need not be optimal.  For comparison, if a moment
vector $y$ has known mean covariance $C/n$ and an ellipsoidal radius $q$, a
trace-one PSD witness $X$ has population signal-to-noise ratio
\begin{equation}
 \kappa(X)=\frac{[-\tr(\mathcal L(y)X)]_+}
 {\sqrt{\ell(X)^{\mathsf T}C\ell(X)}},
 \qquad \ell(X)_k=\tr(L_kX).
 \label{eq:snr}
\end{equation}
The robust upper bound becomes negative exactly when
$\sqrt n\,\kappa(X)>q$.  Optimizing $X$ therefore trades localizer negativity
against covariance sensitivity, rather than simply minimizing the observed
eigenvalue.

Mixed witnesses can be strictly better in principle.  For the synthetic family
$M=-\varepsilon I_r$ with isotropic Frobenius uncertainty, every rank-one
projector has uncertainty norm one, whereas $X=I_r/r$ has norm $r^{-1/2}$.
The sample threshold differs by a factor $r$.  In our ANNNI pilots, however,
many optimized ellipsoidal witnesses above threshold are numerically rank one.
The observed advantage there is primarily \emph{covariance-aware direction
selection} relative to the plug-in minimum eigenvector, not evidence that a
high-rank mixture is always required.

\section{Interacting ANNNI experiment}

We use the periodic axial next-nearest-neighbour Ising chain
\begin{equation}
 H=-J_1\sum_i Z_iZ_{i+1}-J_2\sum_i Z_iZ_{i+2}-h_x\sum_iX_i
 \label{eq:annni}
\end{equation}
with $J_1=1$, $J_2=0.37$, $h_x=2.2$, and $\tau=0.2$.  This frustrated,
nonintegrable finite chain~\cite{Kriel2014} is not being used to claim a new phase boundary; it
is a controlled interacting test bed where exact sparse diagonalization and
Krylov propagation provide ground truth.  Ground-state-subtracted central-site
$X$ and $Z$ probes form a two-channel correlator.  The global spin-flip parity
separates their visible sectors.

At degree $N=2$, the primitive joint sketch dimension is $d=12$, the flattened
quadratic moment dimension is 18, and the localizer has size $6\times6$.
Lengths $L=6,8,10,12,14,16$ are tested at the exact visible edge and at gaps
inflated by 15\% and 35\%.  For each cell we generate independent sketch
ensembles at $n=500,1000,2000,5000$.  The production seed, full raw summaries,
and SHA-256 hashes are stored in the repository.

\subsection{Baselines and controls}

We compare four logically different objects.
\begin{enumerate}
 \item The Wishart--Loewner dual certificate of Proposition~\ref{prop:dual}, with unknown
 covariance and finite-sample level control.
 \item A covariance-known multivariate-Chebyshev ellipsoid, useful as a
 controlled oracle reference but not a deployable unknown-covariance method.
 \item A plug-in minimum-eigenvector and a fixed Chebyshev polynomial filter.
 These expose the value of covariance-aware direction selection.
 \item A generalized Ritz/GEVP edge estimate.  It is a point estimator, not a
 confidence certificate, and is reported separately rather than granted the
 same logical status.
\end{enumerate}

The known-covariance pilot showed a sharp sample transition: for the 35\%
inflated gap, the covariance-aware robust test reached essentially full power
near $n=1000$ across $L\le14$, while the plug-in eigenvector was usually still
silent until $n=5000$.  At the true edge the robust methods made no rejection
in the medium grid.  These observations motivated the covariance-unknown
production campaign.

\subsection{Production power and moment-degree transition}

The main unknown-covariance campaign comprises 7,200 independently generated
SDPs.  It made zero rejections in 2,400 null ensembles and 2,400 rejections in
2,400 ensembles at 35\% inflation.  For the harder 15\% alternative, the
finite-size transition is resolved in Table~\ref{tab:mainpower}.  Each entry is
the number of rejections among 100 ensembles; Wilson intervals are plotted in
Figure~\ref{fig:powerdegree}.  Across all 7,200 solves, the largest conservative
stationarity allowance was $3.76\times10^{-13}$ and the smallest reported
eigenvalue of a shifted multiplier was $1.40\times10^{-14}$.

\begin{table}[H]
\centering
\caption{Safe rejections out of 100 for a 15\% inflated visible-gap claim at
moment degree $N=2$.}
\label{tab:mainpower}
\begin{tabular}{c@{\qquad}rrrr}
\toprule
$L$ & $n=500$ & $n=1000$ & $n=2000$ & $n=5000$\\
\midrule
6  &100&100&100&100\\
8  &100&100&100&100\\
10 & 98&100&100&100\\
12 & 51&100&100&100\\
14 &  0& 86&100&100\\
16 &  0&  7& 97&100\\
\bottomrule
\end{tabular}
\end{table}

\begin{figure}[t]
\centering
\includegraphics[width=\textwidth]{paper/figures/power_and_degree.pdf}
\caption{Finite-sample power.  (a) The $N=2$ production campaign with 95\%
Wilson intervals.  (b--d) Independent degree ablation at 15\% inflation.
Increasing $N$ raises the primitive dimension from 8 to 16 and widens the
Loewner band, yet the stronger polynomial localizer dominates: at $L=16$ and
$n=500$, power rises from 0 for $N=1,2$ to 0.90 for $N=3$.}
\label{fig:powerdegree}
\end{figure}

The independent 7,200-SDP degree ablation made zero null rejections in 3,600
ensembles.  Degree one is powerless at $L=12,16$ throughout the tested range;
degree three reaches unit power in every $L=8,12$ cell and at $n\ge1000$ for
$L=16$.  This is not the tautology that adding moments can only help: the
Wishart dimension penalty grows with $N$.  The observed dominance therefore
quantifies a favorable approximation--concentration tradeoff for this model.

\subsection{Conservatism of the analytic band}

In 10,000 standard-Gaussian calibration replicates at $d=12$, the analytic
$\eta$ values were $0.276,0.195,0.138$ for $n=500,1000,2000$, whereas the
empirical 95th percentiles of the required singular-value deviation were
$0.173,0.122,0.0862$.  The analytic width was consistently about 1.60 times
the empirical quantile and covered all 10,000 replicates in each cell.  Thus
the strong production power is not obtained by tuning the band to its Monte
Carlo distribution.

\subsection{Negative calibration result}

The quadratic sketch features obey Isserlis' covariance formula but are
Wishart-like rather than multivariate Gaussian.  At $L=12$, $N=2$, and only
$n=120$ sketches, a nominal 95\% Hotelling ellipsoid covered the exact moment
vector in 83.4\% of 500 simulated ensembles.  Its apparent 97.4\% detection
rate for an inflated gap was therefore invalid as a 95\% certificate.  Coverage
approached nominal only slowly with $n$.  This failed route is retained in the
artifact as a warning against importing an exact Gaussian-mean theorem into a
quadratic-feature problem.

\subsection{Primitive-law misspecification}

The Wishart--Loewner theorem is exact for Gaussian primitive sketches, not for
all zero-mean isotropic laws.  We tested this boundary in 2,400 further SDPs,
using Gaussian, variance-matched Student-$t_5$, and Rademacher primitive
coordinates before applying the same population square root.  The Gaussian
band covered in all 800 in-model ensembles.  Rademacher sketches also happened
to cover in all 800 tested ensembles, but this is empirical conservatism, not a
theorem.  Student-$t_5$ coverage was only 0.780 at $n=500$ and 0.738 at
$n=1000$ after pooling the two independently generated gap-factor cells; its
heavy radial tails do not become licensed merely by covariance matching.

\begin{figure}[t]
\centering
\includegraphics[width=0.92\textwidth]{paper/figures/misspecification_coverage.pdf}
\caption{Coverage stress test for the simultaneous Loewner band, with Wilson
intervals.  Each bar pools 400 ensembles.  Only the Gaussian bars are covered
by Theorem~\ref{thm:wishart}; the other two laws are deliberate negative
controls.}
\label{fig:misspecification}
\end{figure}

\section{Numerical verification protocol}

Ill-conditioned Euclidean moment matrices make raw primal status unreliable.
In a first primal--dual Wishart replay, two of five inflated-gap samples were
labeled \texttt{optimal\_inaccurate} by a primal solver even though their dual
margins were below $-4\times10^{-4}$.  In a separate ten-ensemble $L=8$ smoke
at $n=500$, null dual objectives lay between $1.8\times10^{-10}$ and
$6.5\times10^{-10}$, while 35\%-inflated-gap objectives lay between
$-8.2\times10^{-4}$ and $-4.1\times10^{-4}$.  The maximum stationarity residual
was $2.3\times10^{-9}$.  The first production pass conservatively thresholded
the raw dual margin at $-10^{-7}$ and stored all residual diagnostics.  The
final pass uses the feasibility repair described next.

\paragraph{Feasibility repair.}
Let floating-point multipliers first be projected onto the PSD cone, and let
$R=Y_A-Y_B+\mathcal L_D^*(Y_L)=R_+-R_-$ be their remaining stationarity
residual.  The replacements
\begin{equation}
  \widehat Y_A=Y_A+R_-,\qquad
  \widehat Y_B=Y_B+R_+
  \label{eq:repair}
\end{equation}
remain PSD and satisfy stationarity exactly, up to the arithmetic used for the
eigendecomposition.  Trace renormalization preserves this equality.  In a
20-ensemble repaired smoke, the maximum stationarity residual fell below
$8\times10^{-17}$ and normalization error below $4\times10^{-16}$; null
objectives were positive, whereas all 15\% and 35\% inflated-gap objectives
were negative.  The production threshold is $-10^{-9}$ on the \emph{repaired}
objective, leaving an additional buffer above double-precision roundoff.

For an additional a-posteriori check, let $\widehat R$ be the final residual and
recall that every scaled primal matrix obeys $0\preceq\widetilde Q\preceq B$.
Then
\begin{equation}
  |\ip{\widehat R}{\widetilde Q}|
  \le \|\widehat R\|_*\,\|B\|_{\rm op}.
  \label{eq:residualbound}
\end{equation}
We add this worst-case quantity to the repaired dual objective before making a
decision.  The code also shifts each multiplier by an explicit roundoff-scale
identity so that its reported minimum eigenvalue is positive.  The final rule
therefore tests a conservative upper bound on the dual contradiction margin,
not the raw solver objective.

This gap between scientific margin and numerical residual is encouraging but
does not by itself constitute exact arithmetic.  The release uses three
increasing levels of assurance:
\begin{enumerate}
 \item solve in floating point after diagonal congruence scaling;
 \item independently rebuild the adjoint and residual from exported matrices;
 \item repair PSD and stationarity defects while adding their worst-case effect
 to the contradiction margin;
 \item replay one wider-band representative with exact rational arithmetic.
\end{enumerate}

\subsection{Exact-rational representative witness}
\label{sec:rational}

For $L=8$, $N=2$, $n=500$, and the rational proposed edge
$\theta=3/5$, we replace the analytic deviation by the larger rational value
$\bar\eta=3/10$.  This widening is itself certified by elementary rational
bounds: $\exp(15/4)>40$ via a positive Taylor partial sum,
$\sqrt{12}<7/2$, $\sqrt{2\log40}<11/4$, and
$25/4<3\sqrt5=(3/10)\sqrt{500}$.  Hence infeasibility for
\begin{equation}
 \frac{100}{169}S\preceq K\preceq\frac{100}{49}S
 \label{eq:rationalband}
\end{equation}
implies infeasibility for the narrower 95\% analytic band.

The stored Gaussian coordinates are hexadecimal IEEE-754 values.  The replay
interprets each as its exact dyadic rational, rebuilds $S$ by rational dot
products, and reads rounded dual multipliers with a common denominator
$10^{14}$.  A diagonal shift of $10^{-11}$ makes all three multipliers strictly
positive definite.  Exact fraction LDL decompositions give a minimum pivot
$2.85\times10^{-11}$.  Instead of a floating nuclear norm, the replay uses the
coarser exact inequality
\begin{equation}
 |\ip{R}{Q}|\le d\max_{ij}|R_{ij}|\,\tr B,
 \qquad 0\preceq Q\preceq B.
 \label{eq:rationalresidual}
\end{equation}
After this allowance, the exact rational certificate upper bound is
$-6.47697914332\times10^{-4}$.  The JSON artifact has SHA-256 (line break is
typographical)
\begin{center}\small
\texttt{aa9e6af05079ca823b392d2ee4889928}\\[-2pt]
\texttt{2b87d32839021bd84b6085bd5f8d842d}
\end{center}
and is independently replayed by
\href{https://github.com/lluiseriksson/finite-sample-spectral-certificates/blob/main/verification/verify_rational_witness.py}{\texttt{verify\_rational\_witness.py}}.
This exact statement is conditional on the stored finite-precision sample; the
Gaussian idealization supplies the statistical coverage statement, not the
rational replay itself.

\section{Limitations and interpretation}

\paragraph{Visible versus Hamiltonian gap.}
The probes may have zero overlap with the first excitation.  Multiple symmetry
channels reduce but do not remove this risk.  Every result is consequently
labeled $\Delta_{\rm vis}$ unless a separate overlap theorem is available.

\paragraph{Upper versus lower bounds.}
Rejecting support in $[0,\theta]$ proves that a claimed gap is too large.  A
finite localizer's non-rejection is not a proof that the true support lies below
$\theta$.  Our method therefore complements, rather than replaces, lower-gap
Hamiltonian certificates~\cite{Rao2025}.

\paragraph{Acquisition model.}
The exact confidence band uses independent, known-zero-mean Gaussian sketch
coordinates.  Classical stochastic estimators or randomized numerical linear
algebra can realize this model directly; a laboratory implementation must state
how those coordinates are obtained and account for state-preparation and
measurement noise.  Markov-chain autocorrelation, non-Gaussian primitive
sketches, and unknown centering require different concentration results.

\paragraph{Finite volume.}
The ANNNI study validates statistical behavior at finite $L$.  It neither
establishes a thermodynamic gap nor competes with the convergent bulk hierarchy
of Xu et al.~\cite{Xu2026}.  A future combination of noisy-data confidence sets
with thermodynamic operator-algebra hierarchies is an open direction.

\section{Conclusion}

The main result is a finite-sample bridge from random correlator sketches to an
auditable spectral-edge exclusion.  A single Loewner band absorbs all
time-channel correlations and all adaptive witness selection; a semidefinite
alternative turns rejection into a replayable matrix inequality.  The approach
is intentionally one-sided.  It can falsify an overly large visible gap with
controlled type-I error, but it does not infer a mass gap from silence.  The
interacting-chain campaign is designed to quantify when that rigorous statement
is statistically useful and when conventional point estimates are merely more
optimistic.

\paragraph{Reproducibility.}
Code, theorem audit notes, raw pilot outputs, the Colab production notebook, and
the frozen production artifacts are available at
\url{https://github.com/lluiseriksson/finite-sample-spectral-certificates}.
The three production JSON files have SHA-256 digests recorded in
\texttt{results/production/summary.json}; the continuous-integration workflow
rechecks the localizer adjoint, primal--dual smoke, safe Wishart smoke, and
exact-rational replay.

\paragraph{Acknowledgements.}
OpenAI Codex assisted with code generation, numerical orchestration, literature
triage, and adversarial consistency checks.  The author retains responsibility
for all claims and the final manuscript.

\appendix
\section{Strong alternative from compact interval separation}
\label{app:duality}

We spell out the separation argument behind Proposition~\ref{prop:dual}, since
generic semidefinite systems can exhibit weak infeasibility.  Let
\[
 \mathcal D=\{Q:A\preceq Q\preceq B\},\qquad
 \mathcal M=\{\mathcal L_D(Q):Q\in\mathcal D\}.
\]
When $A\prec B$, $\mathcal D$ is compact, convex, and has a relative-interior
point; hence $\mathcal M$ is compact and convex.  If
$\mathcal M\cap\PSD=\varnothing$, strong separation from the closed PSD cone
provides a nonzero matrix $Y_L\succeq0$ such that
\begin{equation}
 \sup_{Q\in\mathcal D}\ip{Y_L}{\mathcal L_D(Q)}<0.
 \label{eq:separation}
\end{equation}
The separator must be PSD: otherwise its linear functional is unbounded below
or above along a PSD ray, with the orientation chosen in
\eqref{eq:separation}.

The support problem on the left of \eqref{eq:separation} is an SDP with a
strictly feasible interval point $(A+B)/2$.  Its Lagrange dual therefore
attains the same optimum.  Introducing $Y_A,Y_B\succeq0$ for
$Q-A\succeq0$ and $B-Q\succeq0$ makes the coefficient of $Q$ vanish exactly
when
\[
 Y_A-Y_B+\mathcal L_D^*(Y_L)=0,
\]
and leaves the constant
$-\ip{Y_A}{A}+\ip{Y_B}{B}$.  Strict negativity follows from
\eqref{eq:separation}.  This proves the nontrivial direction of the alternative
without invoking a blanket ``SDP Farkas lemma'' in a regime where it might be
false.  Conversely, pairing any feasible $Q$ with such multipliers gives an
immediate contradiction, as in the main text.

If $S$ is singular, then $A$ and $B$ have the same kernel and no ambient strict
midpoint exists.  Because $A\preceq Q\preceq B$ forces
$\ker S\subseteq\ker Q$, one may compress every matrix to
$\operatorname{ran}S$ and repeat the argument there.  The current production
ensembles have $n>d$ and a full-rank population Gram; the explicit range
reduction is nevertheless preferable to adding an undocumented ridge.

\section{Exact robust separation for ellipsoidal moment errors}
\label{app:ellipsoid}

The covariance-known comparator used in the pilot follows a more general
finite-dimensional identity.  Let $y\in\mathbb R^p$ parameterize an affine
localizer
\[
 L(y)=L_0+\sum_{k=1}^p y_kL_k
\]
and let
\[
 \mathcal E=\{\widehat y+Rz:\|z\|_2\le q\}
\]
be any compact ellipsoid.  For a trace-one $X\succeq0$ define
$\ell(X)_k=\tr(L_kX)$ and
\begin{equation}
 U_{\mathcal E}(X)=\tr[L(\widehat y)X]
       +q\|R^{\mathsf T}\ell(X)\|_2.
 \label{eq:ellipsoidupper}
\end{equation}

\begin{theorem}[Exact ellipsoidal incompatibility]
\label{thm:ellipsoid}
The intersection $\mathcal E\cap\{y:L(y)\succeq0\}$ is empty if and only if
\begin{equation}
 \min_{X\succeq0,\ \tr X=1}U_{\mathcal E}(X)<0.
 \label{eq:ellipsoidtest}
\end{equation}
\end{theorem}

\begin{proof}
The image $L(\mathcal E)$ is compact and convex.  If it is disjoint from the
closed PSD cone, strong separation supplies a PSD separator $X$; normalization
to trace one only fixes its nonzero scale.  The support function of the
ellipsoid in that direction is exactly \eqref{eq:ellipsoidupper}.  Conversely,
if \eqref{eq:ellipsoidtest} is negative, then
$\tr[L(y)X]<0$ for every $y\in\mathcal E$, which is impossible when both
$L(y)$ and $X$ are PSD.
\end{proof}

For a singleton ellipsoid, \eqref{eq:ellipsoidtest} reduces to
$\lambda_{\min}(L(y))$ and admits a rank-one eigenvector witness.  Under
uncertainty, however, the optimizer minimizes a sum of signal and sensitivity.
It may be higher rank, and even when it is rank one it need not be the observed
minimum eigenvector.  The same coverage-event argument as in
Corollary~\ref{cor:level} shows that optimizing $X$ on the data does not inflate
type-I error when $\mathcal E$ is a simultaneous confidence set.

\section{Quadratic-feature covariance and the Hotelling control}
\label{app:isserlis}

For completeness, Isserlis' identity gives the exact covariance of one
quadratic sketch replicate.  With joint covariance entries
$K_{ka,\ell c}=\mathbb E[z_{k,a}z_{\ell,c}]$,
\begin{align}
 &\operatorname{Cov}(z_{k,a}z_{k,b},z_{\ell,c}z_{\ell,d})\\
 &\qquad=K_{ka,\ell c}K_{kb,\ell d}
          +K_{ka,\ell d}K_{kb,\ell c}.
 \label{eq:isserlis}
\end{align}
For the mean of $n$ independent sketches this matrix is divided by $n$.
If it were known, multivariate Chebyshev would give a distribution-free
ellipsoid with squared radius $r/\alpha$ on its rank-$r$ range.  That oracle
route is exact but conservative and requires the unknown $K$ through
\eqref{eq:isserlis}; it is used only as a controlled comparator.

By contrast, the standard one-sample Hotelling pivot is exact when the
\emph{feature vectors themselves} are multivariate Gaussian.  Products of
Gaussian coordinates violate that hypothesis.  The 83.4\% coverage observed
at nominal 95\% and $n=120$ is therefore not a paradox; it is a direct stress
test of a misapplied calibration.  The Wishart--Loewner route avoids this error
by constructing its pivot before the quadratic map.

\bibliographystyle{unsrt}
\bibliography{references}
\end{document}
